%% =====================================================================
%%  T-23-04  The Gut Verdict
%%  -------------------------------------------------------------------
%%  One idea per file. Core is mandatory. Slots in canonical order.
%%  Max 700 words. No layout commands. No filler. New source material
%%  for this entry goes in sources/B6/T-23-04.tex -- never by
%%  rewriting this file (docs/RULES.md R0.1).
%%  =====================================================================
%% @kind: topic
%% @id: T-23-04
%% @title: The Gut Verdict
%% @subtitle: The confusion you feel is the data --- trust the alarm before the explanation
%% @chapter: c23
%% @order: 40
%% @status: stable
%% @sources: B6, B7
%% @updated: 2026-09-19

\topic{T-23-04}{The Gut Verdict}{The confusion you feel is the data --- trust the alarm before the explanation}

\begin{core}
The most reliable manipulation detector most people own is the one they talk themselves out of. B6's rule: heighten your gut-level sensitivity --- the unease you cannot yet justify is pattern-recognition running ahead of argument, and the manipulator's first target is precisely that signal. The discipline is to treat the alarm as a finding, then do the audit --- never to explain the alarm away before looking.
\end{core}
\begin{mechanism}
B6's patients all reported the same instrument readings --- distrust they could not name a reason for --- and had all been taught by the manipulator to override them: too suspicious, too sensitive, imagining things. The mechanism is that covert tactics are engineered to be deniable at the level of conscious analysis (the veiled threat, the innocence look, the vague answer), while leaving a legible residue at the level of felt discomfort; the disguise defeats argument but not alarm. That is why B6's training order is exact: commit the tactic list to memory, attend closely to conduct --- track what the words are doing, not only what they say --- and when the gut speaks, audit before discounting. The named vulnerabilities are the override's codebook: naivete (refusing to believe anyone is as cunning as your gut reports), over-conscientiousness (pre-crediting their side), low self-confidence (doubting your right to the judgement), over-intellectualisation (believing that if you understood the reasons, the reasons would excuse the behaviour). Each turns a working detector off at the moment of reading.
B7's interrogation practice confirms the alarm and disciplines it: unexamined distrust is only global assessment --- guessing --- so the gut selects the moment for a question and a scored window does the reading; its obstacle list (nobody would lie to me; calling a liar is worse than lying) is the override's social reinforcement.
\end{mechanism}
\begin{conditions}
Strongest on first contact and in long relationships gone ambient --- the two places analysis is weakest (no history yet; habituation). It weakens wherever the alarm has a cheap alternative explanation you have verified (stress, hunger, their documented bad week) --- the audit exists to find those. It fails entirely when run as pure paranoia: the gut can also be trained on false patterns, which is why the finding must go to audit, not straight to verdict.
\end{conditions}
\begin{application}
  \item Log the alarm when it fires: who, what happened, what exactly feels off. Patterns need records; feelings alone gaslight themselves.
  \item Listen for, not to: track what the words are doing (diverting, vague, guilt-loading) alongside what they are saying (T-06-02).
  \item Run the vulnerability checklist against your own discount: am I overriding this because of naivete, over-conscientiousness, self-doubt, or the need to understand first?
  \item Give the finding one honest audit before any override. Not after; before.
  \item When the audit confirms, act on behaviour --- distance, records, refusal --- not on accusation (T-24-06).
\end{application}
\begin{feedback}
  \item Working: the log shows the alarm clustering around one person's requests. Verdict follows evidence; the feeling was a flag, not a fact.
  \item Working: you catch yourself composing their defence and stop. The override code is losing its grip.
  \item Failing: you feel permanently suspicious of everyone. The instrument is over-gained; recalibrate with the audit step, which you skipped.
  \item Failing: the alarm fires and you talk yourself out of it in under a minute. That speed is the vulnerability, not the verdict.
\end{feedback}
\begin{failure}
The gut verdict unexamined becomes prejudice with a body: it convicts accents, mannerisms and difference on the same circuitry that convicts real wolves. And over-trusting it isolates --- a person who audits nothing eventually trusts no one, and the manipulator's final victory is a target who can no longer tell ally from operator because they stopped running evidence at all.
\end{failure}
\begin{countermeasures}
  \item In your own dealings, honour the other party's unease when they raise it --- especially when your hands are clean. Dismissing it is what the wolf does (T-23-02).
  \item Audit your dismissals: how often have you told someone \emph{you're imagining things} about their own discomfort? The phrase is the override, spoken aloud.
  \item Keep your instincts calibrated with rest and records. The detector degrades exactly as the log goes dark.
  \item If someone's charm repeatedly leaves you disoriented rather than delighted, treat the disorientation as the invoice arriving early (T-03-03).
\end{countermeasures}

\sources{B6, B7}
\seesources
\seealso{T-03-03, T-23-01, T-06-02}
